\chapter{Gaussian Tests for the Signed Gamma Form}
\label{v4-arith:w7-d:chapter}

A Gaussian decays fast enough for every term of the explicit formula,
although it is not compactly supported. This extension requires
exponential control in the logarithmic variable. Ordinary Schwartz
continuity cannot control the uncentered critical prime measure.

\section{An exponential test space}
\begin{definition}[rapid exponential decay]
Let $\mathscr G$ consist of all smooth functions $f$ such that
\[
 p_{a,N,j}(f)=\sup_{t\in\mathbb R}
 e^{a|t|}(1+|t|)^N|f^{(j)}(t)|<\infty
 \quad(a,N,j\in\mathbb N).
\]
Give $\mathscr G$ the topology of these seminorms. Use
$h_f(z)=\int f(t)e^{izt}dt$ and $f^*(t)=\overline{f(-t)}$.
\end{definition}

\begin{proposition}[extension of the explicit formula]
\label{v4-arith:sc:gaussian-extension}
The space $\mathscr G$ is a convolution algebra with continuous
involution. Compact smooth functions are dense in $\mathscr G$.
The explicit formula, its signed pairing, and both polar moments
extend continuously to $\mathscr G$.
In particular, for $f,g\in\mathscr G$,
\[
 W(f*g^*)=\sum_\rho h_f(z_\rho)\overline{h_g(\bar z_\rho)}
 =P(f*g^*)-W_\infty(f*g^*)-W_{\mathrm{fin}}(f*g^*).
\]
The zero and prime sums are absolutely convergent. Positivity on
all squares in $\mathscr G$ is equivalent to positivity on all
squares in $\mathscr A$.
\end{proposition}
\begin{proof}
The inequality $|t|\le|v|+|t-v|$, together with the corresponding
polynomial weight inequality, bounds every convolution seminorm
by a weighted supremum seminorm times a weighted integral seminorm.
An additional exponential weight bounds that integral. This proves
continuity of convolution. Reflection proves continuity of $*$.
For a compact smooth cutoff $\chi=1$ near zero, put
$f_R(t)=\chi(t/R)f(t)$. Leibniz's rule and one extra exponential
weight give $p_{a,N,j}(f-f_R)\to0$ for every seminorm.

Integration by parts gives rapid decay of $h_f$ in each fixed
horizontal strip, continuously in these seminorms. The zero count
then bounds the zero sum. The polar evaluations are continuous.
The gamma density has logarithmic growth, so its integral is
continuous. Finally, for example,
\[
 \sum_{n\ge2}\frac{\Lambda(n)}{\sqrt n}|f(\pm\log n)|
 \le p_{2,0,0}(f)\sum_{n\ge2}\frac{\log n}{n^{5/2}}.
\]
The compact explicit formula therefore passes to the limit $f_R$.
Convolution continuity gives the paired identity and passage of
positivity. The reverse positivity implication uses
$\mathscr A\subset\mathscr G$.
\end{proof}

\section{The normalized scalar integral}
\begin{lemma}[Gaussian Fourier transform]
For $\lambda>0$ set $f_\lambda(t)=e^{-\lambda t^2/2}$. Then
\[
 h_{f_\lambda}(u)=\sqrt{\frac{2\pi}{\lambda}}
 e^{-u^2/(2\lambda)},\qquad
 (f_\lambda*f_\lambda)(t)=\sqrt{\frac\pi\lambda}
 e^{-\lambda t^2/4}.
\]
Consequently, with $k(u)=\log\pi-\Re\psi_\Gamma(1/4+iu/2)$,
\[
 W_\infty(f_\lambda*f_\lambda)
 =\sqrt{\frac\pi\lambda}\,H(\lambda),\qquad
 H(\lambda)=\frac1{\sqrt{\pi\lambda}}
             \int e^{-u^2/\lambda}k(u)du.
\]
\end{lemma}
\begin{proof}
Completing the square gives both Gaussian integrals.
Insert $|h_{f_\lambda}|^2=(2\pi/\lambda)e^{-u^2/\lambda}$
in the gamma pairing with its coefficient $1/(2\pi)$.
\end{proof}

\begin{theorem}[convergent error-function series]
\label{v4-arith:sc:erfc}
Let $c_n=n+1/4$ and
$\operatorname{erfc}(x)=\frac2{\sqrt\pi}\int_x^\infty e^{-v^2}dv$.
For every $\lambda>0$,
\[
 H(\lambda)=k(0)-\sum_{n=0}^\infty
 \left[\frac1{c_n}-\frac{2\sqrt\pi}{\sqrt\lambda}
 e^{4c_n^2/\lambda}\operatorname{erfc}(2c_n/\sqrt\lambda)\right].
\]
Each bracket is nonnegative and at most $\lambda/(8c_n^3)$.
The series converges uniformly on compact $\lambda$-intervals.
\end{theorem}
\begin{proof}
The digamma series gives
$k(0)-k(u)=\sum_n (u^2/4)/(c_n(c_n^2+u^2/4))$.
Tonelli's theorem exchanges this nonnegative sum and the Gaussian
integral. For $a>0$,
\[
 \int_{\mathbb R}\frac{e^{-u^2/\lambda}}{u^2+a^2}du
 =\sqrt\pi\int_0^\infty\frac{e^{-a^2r}}{\sqrt{r+1/\lambda}}dr
 =\frac\pi a e^{a^2/\lambda}\operatorname{erfc}(a/\sqrt\lambda).
\]
The first equality follows by integrating
$\int_0^\infty e^{-r(u^2+a^2)}dr$ in either order.
The substitution $v=a\sqrt{r+1/\lambda}$ gives the second.
Set $a=2c_n$ and use
$\int e^{-u^2/\lambda}u^2du=\sqrt\pi\lambda^{3/2}/2$.
The estimate on each bracket follows by replacing
$c_n^2+u^2/4$ by $c_n^2$. It proves uniform convergence.
\end{proof}

\begin{theorem}[the exact Gaussian sign range]
\label{v4-arith:w7-d:thm:gauss-ATP}
The continuous function $H$ is strictly decreasing on $(0,\infty)$,
with limits $k(0)>0$ at zero and $-\infty$ at infinity.
There is a unique $\lambda_c>0$ with $H(\lambda_c)=0$, and
\[
 W_\infty(f_\lambda*f_\lambda)\ge0
 \quad\Longleftrightarrow\quad 0<\lambda\le\lambda_c.
\]
\end{theorem}
\begin{proof}
After $u=\sqrt\lambda v$,
$H(\lambda)=\pi^{-1/2}\int e^{-v^2}k(\sqrt\lambda v)dv$.
Strict radial decrease of $k$ gives strict decrease of this integral.
Logarithmic growth gives a common integrable majorant on each compact
$\lambda$-interval and on $0<\lambda\le1$.
Thus $H$ is continuous and tends to $k(0)$ at zero.
For each $v\ne0$, the nonnegative quantity
$k(0)-k(\sqrt\lambda v)$ increases to infinity.
Monotone convergence gives the limit at infinity.
The intermediate value theorem completes the proof.
\end{proof}

\section{The arithmetic inequality on Gaussian squares}
For these Gaussian tests the polar and prime contributions are
\[
 P(f_\lambda*f_\lambda)=\frac{4\pi}{\lambda}e^{1/(4\lambda)},
 \qquad
 W_{\mathrm{fin}}(f_\lambda*f_\lambda)
 =2\sqrt{\frac\pi\lambda}\sum_{n\ge2}
   \frac{\Lambda(n)}{\sqrt n}e^{-\lambda(\log n)^2/4}.
\]
These expressions converge for every $\lambda>0$ and give the
complete value of $W(f_\lambda*f_\lambda)$ after subtracting the
prime and gamma terms from the polar term.
The gamma sign theorem alone does not determine that value.

\begin{proposition}[a single Gaussian parameter does not test a quadratic form]
For a Hermitian form, nonnegative values on a spanning collection
of vectors do not imply nonnegative values on its linear span.
In particular, Gaussian cross terms require their own comparison.
\end{proposition}
\begin{proof}
The Hermitian matrix
$\left(\begin{smallmatrix}1&2\\2&1\end{smallmatrix}\right)$
is positive on each coordinate vector but has value $-2$ on $(1,-1)$.
More generally, $q(f+g,f+g)=q(f,f)+q(g,g)+2\Re q(f,g)$.
The signed arithmetic representation determines these cross terms
on $\mathscr G$, but positivity of individual Gaussian squares
does not bound them.
\end{proof}
